\documentclass[10pt,reqno]{amsart}

\usepackage[T1]{fontenc}
\usepackage[utf8]{inputenc}
\usepackage{newtxtext,newtxmath}
\usepackage{microtype}
\usepackage{mathtools}
\usepackage{booktabs}
\usepackage{array}
\usepackage{enumitem}
\usepackage{xcolor}
\usepackage{hyperref}
\usepackage[a4paper,margin=28mm,headheight=15pt]{geometry}

\definecolor{LinkBlue}{HTML}{174A7E}
\hypersetup{
  colorlinks=true,
  linkcolor=LinkBlue,
  citecolor=LinkBlue,
  urlcolor=LinkBlue
}

\setlength{\parskip}{2pt}
\setlist[itemize]{leftmargin=6mm,itemsep=1.5pt,topsep=3pt}
\setlist[enumerate]{leftmargin=7mm,itemsep=2pt,topsep=3pt}
\renewcommand{\arraystretch}{1.15}
\emergencystretch=2em

\newtheorem{theorem}{Theorem}[section]
\newtheorem{proposition}[theorem]{Proposition}
\newtheorem{lemma}[theorem]{Lemma}
\newtheorem{corollary}[theorem]{Corollary}
\theoremstyle{definition}
\newtheorem{definition}[theorem]{Definition}
\newtheorem{remark}[theorem]{Remark}

\title[Dispersion Scale Barriers for Representing Integers as $a p^2 + b$ (Erd\H{o}s \#676)]
{Baier--Zhao Dispersion Scale Barriers for Quadratic-Prime Representations $n = a p^2 + b$}
\author{Argus Open Research Project}
\date{September 5, 2026}
\subjclass[2020]{Primary 11P32; Secondary 11N36}
\begin{document}
\begin{abstract}
Erd\H{o}s Problem~\#676 asks whether every sufficiently large integer $n$ can be written in the form $n = a p^2 + b$ with $p$ prime, $a\ge 1$, and $0\le b < p$. We prove that on the signed covariance dispersion route at exponent $\kappa < 1/2$, the Baier--Zhao large-sieve bound for primes in quadratic progressions supplies only $Y^{34+o(1)}\log Y$ against the required $Y^2$ saving, excluding that dispersion interface without resolving the conjecture.
\end{abstract}
\maketitle
\section{Main Scale Analysis}
\begin{proposition}
At remainder threshold $b < p^\kappa$ with $\kappa < 1/2$, the off-diagonal signed covariance sum over moduli $a\le X/Y^2$ requires cancellation below $X^2/Y$, whereas the unconditional Baier--Zhao quadratic large-sieve estimate yields an upper bound of order $Y^{34+o(1)}\log Y$, leaving an unconditional exponent gap on that analytic route.
\end{proposition}
\end{document}
